% Figure from MATH 452 notes, section 16. Compile with the notes preamble.
\begin{tikzpicture}[scale=1.0]
\draw[thick, figB!70, fill=figB!8, name path=bd,
  postaction={decorate},
  decoration={markings,
    mark=at position 0.10 with {\arrow[figB, scale=1.2]{Stealth}},
    mark=at position 0.74 with {\arrow[figB, scale=1.2]{Stealth}}}]
  plot[smooth cycle, tension=0.85] coordinates {(0,0) (2.4,-0.4) (3.6,0.7) (3.1,2.1) (1.4,2.5) (-0.5,1.3)};
\node at (1.2,1.5) {$D$};
% boundary point P on the right side, tangent direction from a nearby point Q further along
\path[name path=h1] (2.6,1.25) -- (4.5,1.25);
\path[name path=h2] (2.6,1.33) -- (4.5,1.33);
\path[name intersections={of=bd and h1, by=P}];
\path[name intersections={of=bd and h2, by=Q}];
\coordinate (T) at ($(P)!1.0cm!(Q)$);
\coordinate (No) at ($(P)!0.9cm!-90:(Q)$);
\coordinate (Ni) at ($(P)!0.9cm!90:(Q)$);
\draw[-stealth, thick, figB] (P) -- (T) node[anchor=south, inner sep=2pt, font=\scriptsize] {$\mathbf{T} = (x', y')$};
\draw[-stealth, thick, figR] (P) -- (No) node[anchor=west, inner sep=2pt, font=\scriptsize] {$\mathbf{n}_{\text{out}} = (y', -x')$};
\draw[-stealth, thick, black!55, dashed] (P) -- (Ni) node[anchor=north, xshift=-10pt, inner sep=3pt, font=\scriptsize, black!70] {$\mathbf{n}_{\text{in}} = (-y', x')$};
\filldraw[black] (P) circle (1.3pt);
% right-angle mark between T and n_out
\coordinate (A) at ($(P)!0.16cm!(T)$);
\coordinate (B) at ($(P)!0.16cm!(No)$);
\draw[black!45, thin] (A) -- ($(A)+(B)-(P)$) -- (B);
\node[figB!80!black, font=\scriptsize, anchor=north] at (1.2,-0.45) {$\gamma = \partial D$, counterclockwise};
\end{tikzpicture}%
